% =====================================================================================
%  Chapter 8 — Discussion
%  Thesis: Comparison of Control Strategies for Interaction-Force Aware Multirotor Teams
%
%  DRAFT SKELETON, 27 September 2026.
%  Sources: Section~\ref{sec:intro-rq}, Section~\ref{sec:intro-contrib}, Section~\ref{sec:setup-validity},
%           Chapters~\ref{ch:results2}--\ref{ch:results3}.
%
%  HONESTY NOTE. Interpretation of results is deferred until C.4/C.3-robot data exist. Each
%  subsection states what evidence would answer the question and what a negative result would mean,
%  following Chapter~\ref{ch:intro}'s falsifiability framing.
%
%  Citation baseline: 2026-08-27 (see docs/thesis/CITATION_BASELINE.md).
% =====================================================================================

\chapter{Discussion}
\label{ch:discussion}

Chapters~\ref{ch:results2} and~\ref{ch:results3} report measurements; this chapter interprets them
against the research questions in Section~\ref{sec:intro-rq}, checks the intended contributions in
Section~\ref{sec:intro-contrib}, and revisits the threats to validity flagged in
Section~\ref{sec:setup-validity}. A negative answer to any question is treated as a result, not as
a failed experiment.

% -------------------------------------------------------------------------------------
\section{RQ1 --- Reactive versus predictive compensation under identical conditions}
\label{sec:disc-rq1}

RQ1 asks, under identical hardware, trajectories, formations and separations, how reactive and
predictive interaction-force compensation compare to the uncompensated baseline in trajectory-tracking
accuracy and residual-force rejection. Evidence comes from Chapter~\ref{sec:res2-tracking} and
residual magnitudes on the same flights. If compensating strategies do not separate from one another
but all beat Strategy~0, the practical conclusion shifts toward ``compensation matters; family
matters less,'' and RQ4 dominates the engineering decision. If they separate strongly, this chapter
must explain whether the gap is robust across repeats or driven by a single scenario class.

\todo[color=red!25]{PLACEHOLDER: Interpretation paragraph once C.4 tables exist --- ranking S0--3,
residual rejection, statistical spread across $n$.}

% -------------------------------------------------------------------------------------
\section{RQ2 --- Regime dependence}
\label{sec:disc-rq2}

RQ2 asks whether the relative advantage of reactive versus predictive compensation depends on
separation, relative velocity, and quasi-static versus transient interaction. The scenario library
was designed to isolate these regimes (Section~\ref{sec:res2-summary}); a crossover supports
conditional guidance (``use predictive in static stacks, reactive in fast crossings''). If no crossover
appears, the discussion must say so and attribute flat rankings to measurement floor, tuning, or
insufficient excitation rather than silently claiming regime independence.

\todo[color=red!25]{PLACEHOLDER: Scenario-class breakdown referencing Table~\ref{tab:res2-when-which}
and train/test splits across classes.}

% -------------------------------------------------------------------------------------
\section{RQ3 --- Two-robot training on three-robot formations}
\label{sec:disc-rq3}

RQ3 asks how large the transfer penalty is for a summed per-neighbour model trained on two-vehicle
data when used on three-vehicle formations, and whether that penalty is small enough to keep the
summed architecture. The question presupposes non-superposition; the contribution is quantitative.
Chapter~\ref{sec:res3-transfer} supplies the numbers; here they are interpreted against deployment
cost (single weight set, permutation invariance) versus retraining or per-formation models.

\todo[color=red!25]{PLACEHOLDER: Interpret open-loop and (if available) closed-loop three-robot
penalty; tie to Section~\ref{sec:res3-limitation} if only logs exist.}

% -------------------------------------------------------------------------------------
\section{RQ4 --- Cost versus performance}
\label{sec:disc-rq4}

RQ4 asks, for each compared strategy, what sensing, onboard compute and training data are required,
and how those costs sit next to tracking and rejection. Chapter~\ref{sec:res2-effort} and the
predictive-path subsection supply performance-side numbers; cost includes RPM sourcing, SD logging,
upload vs flash-resident weights, tuning time per strategy (Table~\ref{tab:setup-threats}), and
C.1 collection burden for Strategy~2/3.

\todo[color=red!25]{PLACEHOLDER: Cost--performance synthesis table (qualitative + quantitative
where measured); no invented rankings.}

% -------------------------------------------------------------------------------------
\section{Contributions revisited}
\label{sec:disc-contrib}

Section~\ref{sec:intro-contrib} listed three intended contributions: a controlled comparison,
a quantification of the summed-model transfer penalty, and an engineering account of cost alongside
performance. The apparatus (scenario library, residual path, training pipeline) was explicitly
excluded as a scientific claim. For each item below, this section will state whether the C.4/C.3
campaign achieved it, partially achieved it, or could not be tested, without upgrading apparatus
description into a result.

\todo[color=red!25]{PLACEHOLDER: Contribution (1) --- controlled comparison achieved? cite
Chapter~\ref{ch:results2} scope and null scenarios.}

\todo[color=red!25]{PLACEHOLDER: Contribution (2) --- transfer penalty quantified? cite
Chapter~\ref{ch:results3}.}

\todo[color=red!25]{PLACEHOLDER: Contribution (3) --- cost account complete? cite RQ4 subsection.}

% -------------------------------------------------------------------------------------
\section{Threats to validity, revisited}
\label{sec:disc-validity}

Chapter~\ref{sec:setup-validity} and Table~\ref{tab:setup-threats} listed threats and mitigations
at design time. After flights, the discussion closes the loop: which mitigations held in practice,
which residual risks materialised, and which limitations remain unavoidable.

\paragraph{Tuning and fairness.} Per-strategy tuning followed by frozen gains is necessary for RQ1;
the discussion will report whether any comparison flight required ad-hoc gain changes and, if so,
exclude or flag those cells.

\paragraph{Measurement floor and null scenarios.} C1--C3 bound how small a reported improvement can
be before it is indistinguishable from noise; revisiting this prevents over-interpreting marginal
RMSE gaps.

\paragraph{Single airframe and session effects.} These were not mitigated by design; results are
scoped to the Crazyflie brushless configuration flown. Interleaving reduces but does not remove
session effects; repeat variance in Chapter~\ref{ch:results2} quantifies what remains.

\paragraph{Simulation.} No thesis claim rests on simulation (Section~\ref{sec:setup-validity});
any sim-only diagnosis (e.g.\ backend choice for SIL) is mentioned only as implementation context,
not as validation of a strategy.

\todo[color=red!25]{PLACEHOLDER: Short paragraph per threat row in Table~\ref{tab:setup-threats} ---
``after data: mitigated / observed / not testable'' once C.4 complete.}
